% !TEX encoding = IsoLatin

%%%%%%%%%%%%%%%%%%%%%%%%%%%%%%%%%%%%%%%%%%%%%%%%%%%%%%%%%%%%%%%
%                                                             %
% MINIATURAS MATEM\'ATICAS DE LA GACETA DE LA RSME            %
%                                                             %
%%%%%%%%%%%%%%%%%%%%%%%%%%%%%%%%%%%%%%%%%%%%%%%%%%%%%%%%%%%%%%%

\documentclass[autocontact]{gaceta}

%---------------------
% La Gaceta se har\'a con esto activado,
% pero no hace falta ahora si no quieres:
\usepackage[T1]{fontenc}
\usepackage{lmodern}

% Puedes usar esto (u otra codificaci\'on) si quieres:
\usepackage[latin1]{inputenc}
\usepackage[spanish]{babel}
\decimalpoint
%---------------------

%---------------------
% Esto ya lo carga el estilo de La Gaceta:
%\usepackage{amsmath, amsthm, amssymb}
%\usepackage{url} % <-- Para p\'aginas web o similar: \url{...}
%\usepackage{graphicx}
%---------------------

%---------------------
% Carga esto para gr\'aficos si lo necesitas:
% \usepackage{wrapfig}
%---------------------

%---------------------
% Carga esto para gr\'aficos si lo necesitas:
%\usepackage{tikz}
%\usepackage{pgfplots}
%---------------------

% PARA BORRADORES:
%\usepackage{color}
%\newcommand\rojo[1]{\textcolor{red}{#1}}

%---------------------
% <<< ESTO SE AJUSTAR\'A AL EDITAR CADA N\'UMERO DE LA GACETA >>>
%\provisional
\newcounter{pageforminiatura}
\setcounter{pageforminiatura}{212} % Siempre ser\'a p\'agina par
\setcounter{page}{1} % No se usa (dejarlo como 1)
\journame{La Gaceta de la RSME}
\yearofpublication{2017}
\volume{20}
\issuenumber{1}
%---------------------

%---------------------
% Deja esto as\'{\i}:
%\belongstopart{Art\'{\i}culos}
%\belongstopart{Secciones}
%\belongstopart{Noticias de la Sociedad}
%\belongstopart{Actualidad}
\nopart
%---------------------

%---------------------
% Usa este tipo de definiciones si las necesitas:
%\theoremstyle{plain}
%\newtheorem*{Prop}{Proposici\'on}
% \newtheorem*{Teo}{Teo}
%\newtheorem*{theorem}{Teorema}
%---------------------

%---------------------
% Deja esto as\'{\i}:
\supertitle{\large Miniaturas matem\'aticas de \thejourname\\[1mm]
\hfill\footnotesize
Vol.~\thevolume\ (\theyearofpublication),
N\'um.~\theissuenumber, P\'ag.~\thepageforminiatura\theprovisional}
%
\noeditor
%---------------------

\let\contacturl\relax

\newcommand{\sgn}{\operatorname{sgn}}

%---------------------
%% ALTERANTIVAS PARA EL T�TULO:
%\title{Una funci�n constante en un intervalo}
\title{Una funci�n que, sin parecerlo, es constante en un intervalo}
%\author{} % Poner \author{} o no poner nada, da igual.
% Para ahorrar espacio, el autor solo se pone en la firma:
\contact{�scar Ciaurri y Juan Luis Varona,
% Dpto. de Matem�ticas y Computaci�n, 
Universidad de La Rioja}
{\texttt{oscar.ciaurri@unirioja.es}, \texttt{jvarona@unirioja.es}}{}
%---------------------


%%%%%%%%%%%%%%%%%%%%%%%%%%%%%%%%%%%%%%%%%%%%%%%%%%%%%%%%%%%%%%%%%%%%%%%
\begin{document}
%%%%%%%%%%%%%%%%%%%%%%%%%%%%%%%%%%%%%%%%%%%%%%%%%%%%%%%%%%%%%%%%%%%%%%%

%%%%%%%%%%%%%%%%%%%%%%%%%%%%%%%%%%%%%%%%%%%%%%%%%%%%%%%%%%%%%%%%%%%%%%%
% Deja esto as\'{\i}:
\previousvspace{-1.8cm}
\makesupertitle
\previousvspace{-0.25cm}
\maketitle
\vspace*{-0.4cm}
\thispagestyle{empty} % <<<<< Importante
%%%%%%%%%%%%%%%%%%%%%%%%%%%%%%%%%%%%%%%%%%%%%%%%%%%%%%%%%%%%%%%%%%%%%%%

%% A PARTIR DE AQU\'I HAY LIBERTAD PARA EL AUTOR:

En los primeros cursos de An�lisis Matem�tico, resulta muy �til disponer
de funciones con definiciones sencillas que tengan comportamientos
extra�os e insospechados a primera vista, de modo que 
los alumnos se convenzan de la necesidad de introducir el rigor.
Aqu� vamos a ver que la funci�n
\[
F(a) = \int_{-1}^1 \frac{\log(x-a)^2}{\sqrt{1-x^2}}\, dx
\]
es constante (igual a $-2\pi\log2$) cuando $a \in [-1,1]$, pero no lo es cuando $|a|>1$.

Comencemos suponiendo $|a|\le 1$. 
Usando $F(-a)=F(a)$ y el cambio de variable $x=\cos t$, se tiene 
\[
2F(a) = \int_{-1}^1  \frac{\log(x^2-a^2)^2}{\sqrt{1-x^2}}\, dx
= \int_{0}^\pi \log(\cos^2 t-\cos^2\alpha)^2\, dt,
\]
donde hemos tomado $a=\cos\alpha$. 
Mediante transformaciones de sumas en productos se prueba que
$(\cos^2 a-\cos^2 b)^2 = \sen^2(a+b)\sen^2(a-b)$ y, de esta forma, 
usando que la funci\'on $\sen^2 z$ es $\pi$-peri\'odica, llegamos a
\begin{multline*}
2F(a) = \int_{0}^\pi \log(\sen(t+\alpha)\sen(t-\alpha))^2\, dt
= \int_{-\alpha}^{\pi-\alpha}\log(\sen^2(t+\alpha))\, dt \\
+ \int_{\alpha}^{\pi+\alpha}\log(\sen^2(t-\alpha))\, dt
= 2\int_{0}^{\pi}\log(\sen^2 t)\, dt = 2F(1).
\end{multline*}

Para $|a|>1$, derivando bajo el signo integral y aplicando el cambio de variable $x=(1-u^2)/(1+u^2)$, 
se tiene 
\[
F'(a) = 2\int_{-1}^1 \frac{dx}{(a-x)\sqrt{1-x^2}}
= 4\int_{0}^\infty \frac{du}{a-1+(a+1)u^2} 
= \frac{2 \sgn(a) \pi}{\sqrt{a^2-1}},
\]
as� que $F(a) = 2\pi\log(|a|+\sqrt{a^2-1})+C$. 
Por continuidad de $F$, $C=F(1)$.
% Por ser $F$ es una funci\'on continua, es claro que $C=F(1)$. 

% Los c\'alculos que conducen a $F(1)$ son comunes en muchos cursos de 
% an\'alisis b\'asico, pero damos una pincelada del proceso. 
Damos una pincelada de los c\'alculos que conducen a $F(1)$.
De las identidades
\[
\int_{\pi/2}^\pi \log(\sen t)\, dt
= \int_0^{\pi/2} \log(\sen t)\, dt
\quad\text{ y }\quad
\int_{0}^{\pi/2} \log(\sen t)\, dt
= \int_0^{\pi/2} \log(\cos t)\, dt
\] 
se obtiene 
\[
F(1) = 2\left(\int_{0}^{\pi/2} \log(\sen t)\, dt
+ \int_0^{\pi/2} \log(\cos t)\, dt\right)
= 2\int_{0}^{\pi/2} \log(\sen(2t)/2)\, dt,
\]
luego $F(1) = -\pi\log 2+F(1)/2$ y $F(1)=-2\pi\log2$.

\enlargethispage{1.5mm}

\medskip

%%%%%%%%%%%%%%%%%%%%%%%%%%%%%%%%%%%%%%%%%%%%%%%%%%%%%%%%%%%%%%%%%%%%%%%
\end{document}
%%%%%%%%%%%%%%%%%%%%%%%%%%%%%%%%%%%%%%%%%%%%%%%%%%%%%%%%%%%%%%%%%%%%%%%
